% !TEX root = ../Main.tex
\chapter{电表内阻对实验的影响}

用电压表、电流表测电阻，依据的是电阻的定义式
\[
R=\frac{U}{I}.
\]
测量值准不准，取决于电压表、电流表怎么接——因为电表本身有内阻（这正是上一章说的系统误差）。下面比较两种接法。

\section{电流表外接（适合测小电阻）}

电压表直接并联在灯（待测电阻）两端，电流表接在外侧。

\begin{center}
\begin{circuitikz}[scale=0.9, american]
    \draw (0,0) to[battery1] (0,2);
    \draw (0,2) -- (1,2);
    \draw (1,2) to[lamp, l=待测] (3,2);
    \draw (1,2) -- (1,3) to[voltmeter] (3,3) -- (3,2);
    \draw (3,2) to[ammeter] (5,2);
    \draw (5,2) -- (5,0) -- (0,0);
\end{circuitikz}
\end{center}

\textbf{电压表示数没问题}（它测的就是灯两端电压）。但\textbf{电压表会分流}：电流表测到的是流过灯和流过电压表的电流之和，比真正流过灯的电流\textbf{偏大}。于是
\[
R_{\text{测}}=\frac{U}{I_{\text{测}}}<R_{\text{真}},
\]
测量值\textbf{小于真实值}。

\rmk{待测电阻越小，流过它的电流越大，电压表那点分流就越微不足道。所以电流表外接\textbf{测"小电阻"误差更小}。}

\section{电流表内接（适合测大电阻）}

电流表与灯串联，电压表测"灯 $+$ 电流表"整体两端的电压。

\begin{center}
\begin{circuitikz}[scale=0.9, american]
    \draw (0,0) to[battery1] (0,2);
    \draw (0,2) -- (1,2);
    \draw (1,2) to[ammeter] (3,2);
    \draw (3,2) to[lamp, l=待测] (5,2);
    \draw (1,2) -- (1,3) to[voltmeter] (5,3) -- (5,2);
    \draw (5,2) -- (5,0) -- (0,0);
\end{circuitikz}
\end{center}

\textbf{电流表示数没问题}（它测的就是流过灯的电流）。但\textbf{电流表会分压}：电压表测到的是灯和电流表两端电压之和，比灯两端真正的电压\textbf{偏大}。于是
\[
R_{\text{测}}=\frac{U_{\text{测}}}{I}>R_{\text{真}},
\]
测量值\textbf{大于真实值}。

\rmk{待测电阻越大，灯两端电压越高，电流表那点分压就越微不足道。所以电流表内接\textbf{测"大电阻"误差更小}。}

\section{如何选：看相对偏差}

到底选哪种接法？要比的不是绝对偏差，而是\textbf{相对偏差}。

\state{两种偏差}{
    设真实值 $X_1$、测量值 $X_2$。
    \[
    \text{绝对偏差}\ \Delta X=X_2-X_1,\qquad
    \text{相对偏差}\ \frac{\Delta X}{X_1}.
    \]
    选接法时看相对偏差谁更小：相对偏差小的那种更准。
}

按上面两节的结论：测\textbf{小}电阻用电流表\textbf{外接}（电压表分流的相对影响小），测\textbf{大}电阻用电流表\textbf{内接}（电流表分压的相对影响小）。

\begin{center}
\begin{circuitikz}[scale=0.8, american]
    \node at (-1.6,1) {\large \textcircled{1}};
    \draw (0,0) to[battery1] (0,2);
    \draw (0,2) -- (0.8,2) to[lamp] (2.8,2);
    \draw (0.8,2) -- (0.8,3) to[voltmeter, l=$U_1$] (2.8,3) -- (2.8,2);
    \draw (2.8,2) to[ammeter, l=$I_1$] (4.6,2);
    \draw (4.6,2) -- (4.6,0) -- (0,0);
\end{circuitikz}
\end{center}

\begin{center}
\begin{circuitikz}[scale=0.8, american]
    \node at (-1.6,1) {\large \textcircled{2}};
    \draw (0,0) to[battery1] (0,2);
    \draw (0,2) -- (0.8,2) to[lamp] (2.8,2);
    \draw (2.8,2) to[ammeter, l=$I_2$] (4.6,2);
    \draw (0.8,2) -- (0.8,3) to[voltmeter, l=$U_2$] (4.6,3) -- (4.6,2);
    \draw (4.6,2) -- (4.6,0) -- (0,0);
\end{circuitikz}
\end{center}

接法 \textcircled{1} 是电流表外接（电压表只跨灯）、\textcircled{2} 是电流表内接（电压表跨"灯 $+$ 电流表"），分别用来测小电阻和大电阻。
